% =========================================================
% 12_eigenvalues_eigenvectors.tex
% Vectors and Linear Algebra
% =========================================================

% =========================================================
% SECTION DIVIDER
% =========================================================

\section{Eigenvalues and Eigenvectors}

% =========================================================
% EIGENVECTOR
% =========================================================

\begin{frame}{Eigenvectors}

Let

\[
A:\mathbb{R}^n\rightarrow\mathbb{R}^n
\]

be a linear transformation.

A non-zero vector $\vect{v}$ is an \textbf{eigenvector} of $A$
if

\[
\boxed{
A\vect{v}=\lambda\vect{v}
}
\]

for some scalar $\lambda$.

The scalar $\lambda$ is called the corresponding
\textbf{eigenvalue}.

\vspace{0.4cm}

Thus, an eigenvector is a direction that is preserved by
the transformation, although its magnitude and possibly
its orientation may change.

\end{frame}


% =========================================================
% EIGENVALUE INTERPRETATION
% =========================================================

\begin{frame}{Meaning of an Eigenvalue}

If

\[
A\vect{v}=\lambda\vect{v},
\]

then:

\begin{itemize}
    \item $\lambda>1$ stretches $\vect{v}$.
    \item $0<\lambda<1$ shrinks $\vect{v}$.
    \item $\lambda=1$ leaves $\vect{v}$ unchanged.
    \item $\lambda=0$ maps $\vect{v}$ to $\mathbf{0}$.
    \item $\lambda<0$ reverses the direction.
\end{itemize}

\vspace{0.3cm}

The important property is that

\[
\boxed{
\vect{v}\text{ remains on the same line after transformation.}
}
\]

\end{frame}


% =========================================================
% EXAMPLE
% =========================================================

\begin{frame}{Example: An Eigenvector}

Consider
\vspace{-0.1cm}
\[
A=
\begin{bmatrix}
2&0\\
0&3
\end{bmatrix}
\]

and
\vspace{-0.1cm}
\[
\vect{v}
=
\begin{bmatrix}
1\\0
\end{bmatrix}.
\]

Then
\vspace{-0.1cm}
\[
A\vect{v}
=
\begin{bmatrix}
2&0\\
0&3
\end{bmatrix}
\begin{bmatrix}
1\\0
\end{bmatrix}
=
\begin{bmatrix}
2\\0
\end{bmatrix}.
\]

Therefore,
\vspace{-0.1cm}
\[
\boxed{
A\vect{v}=2\vect{v}.
}
\]

Hence,

\[
\boxed{
\vect{v}
\text{ is an eigenvector with eigenvalue }
\lambda=2.
}
\]

\end{frame}


% =========================================================
% FINDING EIGENVALUES
% =========================================================

\begin{frame}{Finding Eigenvalues}

Starting from
\vspace{-0.1cm}
\[
A\vect{v}=\lambda\vect{v},
\]

we obtain
\vspace{-0.1cm}
\[
A\vect{v}-\lambda\vect{v}
=
\mathbf{0}.
\]

Therefore,
\vspace{-0.1cm}
\[
(A-\lambda I)\vect{v}
=
\mathbf{0}.
\]

For a non-zero solution $\vect{v}$ to exist,
\vspace{-0.1cm}
\[
\boxed{
\det(A-\lambda I)=0.
}
\]

This equation is called the
\vspace{-0.1cm}
\[
\boxed{
\text{characteristic equation}.
}
\]

Its roots are the eigenvalues of $A$.

\end{frame}


% =========================================================
% CHARACTERISTIC POLYNOMIAL
% =========================================================

\begin{frame}{Characteristic Polynomial}

The polynomial

\[
\boxed{
p(\lambda)=\det(A-\lambda I)
}
\]

is called the \textbf{characteristic polynomial}.

The eigenvalues satisfy

\[
\boxed{
p(\lambda)=0.
}
\]

For an $n\times n$ matrix, the characteristic polynomial
has degree $n$.

\vspace{0.4cm}

Therefore, finding eigenvalues involves:

\[
\boxed{
\text{Matrix}
\rightarrow
\text{Characteristic polynomial}
\rightarrow
\text{Roots}
}
\]

\end{frame}


% =========================================================
% 2x2 EIGENVALUE EXAMPLE
% =========================================================

\begin{frame}{Example: Finding Eigenvalues}

Consider
\vspace{-0.1cm}
\[
A=
\begin{bmatrix}
2&1\\
1&2
\end{bmatrix}.
\]

The characteristic equation is
\vspace{-0.1cm}
\[
\det(A-\lambda I)=0.
\]

Thus,
\vspace{-0.1cm}
\[
\det
\begin{bmatrix}
2-\lambda&1\\
1&2-\lambda
\end{bmatrix}
=0.
\]

Therefore,
\vspace{-0.1cm}
\[
(2-\lambda)^2-1=0.
\]

Hence,
\vspace{-0.1cm}
\[
(\lambda-1)(\lambda-3)=0.
\]

The eigenvalues are
\vspace{-0.1cm}
\[
\boxed{
\lambda_1=1,
\qquad
\lambda_2=3.
}
\]

\end{frame}


% =========================================================
% FINDING EIGENVECTORS
% =========================================================

\begin{frame}{Finding Eigenvectors}

Once an eigenvalue $\lambda$ is known, substitute it into
\vspace{-0.1cm}
\[
\boxed{
(A-\lambda I)\vect{v}=\mathbf{0}.
}
\]

The non-zero solutions form the eigenvectors associated
with $\lambda$.

For
\vspace{-0.1cm}
\[
A=
\begin{bmatrix}
2&1\\
1&2
\end{bmatrix},
\]

consider $\lambda=3$:
\vspace{-0.1cm}
\[
A-3I
=
\begin{bmatrix}
-1&1\\
1&-1
\end{bmatrix}.
\]

Solving
\vspace{-0.2cm}
\[
(A-3I)\vect{v}=\mathbf{0}
\]

gives
\vspace{-0.2cm}
\[
\boxed{
\vect{v}
=
t
\begin{bmatrix}
1\\1
\end{bmatrix},
\quad t\neq0.
}
\]

\end{frame}


% =========================================================
% EIGENSPACE
% =========================================================

\begin{frame}{Eigenspace}

For an eigenvalue $\lambda$, the corresponding
\textbf{eigenspace} is
\vspace{-0.1cm}
\[
\boxed{
E_{\lambda}
=
\operatorname{Null}(A-\lambda I).
}
\]

Thus, an eigenspace contains
\vspace{-0.1cm}
\[
\text{the zero vector}
\]

and all eigenvectors associated with $\lambda$.

For example, if
\vspace{-0.1cm}
\[
E_{\lambda}
=
\operatorname{span}
\left\{
\begin{bmatrix}
1\\1
\end{bmatrix}
\right\},
\]

then every non-zero vector of the form
\vspace{-0.1cm}
\[
t
\begin{bmatrix}
1\\1
\end{bmatrix}
\]

is an eigenvector corresponding to $\lambda$.

\end{frame}


% =========================================================
% EIGENVECTORS AND NULL SPACE
% =========================================================

\begin{frame}{Eigenvectors and the Null Space}

The eigenvector equation can be written as

\[
(A-\lambda I)\vect{v}
=
\mathbf{0}.
\]

Therefore,

\[
\boxed{
\vect{v}
\in
\operatorname{Null}(A-\lambda I).
}
\]

Hence,

\[
\boxed{
E_{\lambda}
=
\operatorname{Null}(A-\lambda I).
}
\]

This connects eigenvectors directly to the concepts of
null space and homogeneous systems.

\end{frame}


% =========================================================
% DISTINCT EIGENVALUES
% =========================================================

\begin{frame}{Eigenvectors with Distinct Eigenvalues}

Suppose
\vspace{-0.1cm}
\[
A\vect{v}_1=\lambda_1\vect{v}_1
\]

and
\vspace{-0.1cm}
\[
A\vect{v}_2=\lambda_2\vect{v}_2,
\]

where
\vspace{-0.1cm}
\[
\lambda_1\neq\lambda_2.
\]

Then $\vect{v}_1$ and $\vect{v}_2$ are linearly independent.

More generally,

\[
\boxed{
\text{Eigenvectors corresponding to distinct eigenvalues
are linearly independent.}
}
\]

This result is important for constructing an eigenvector basis.

\end{frame}


% =========================================================
% EIGENVECTOR BASIS
% =========================================================

\begin{frame}{Eigenbasis}

Suppose an $n\times n$ matrix has $n$ linearly independent
eigenvectors:

\[
\vect{v}_1,\vect{v}_2,\ldots,\vect{v}_n.
\]

Then these vectors form a basis of $\mathbb{R}^n$.

Construct

\[
\boxed{
P=
\begin{bmatrix}
|&|&&|\\
\vect{v}_1&\vect{v}_2&\cdots&\vect{v}_n\\
|&|&&|
\end{bmatrix}.
}
\]

The eigenvectors provide a coordinate system in which
the transformation becomes particularly simple.

\end{frame}


% =========================================================
% DIAGONAL MATRIX
% =========================================================

\begin{frame}{Diagonalization}

Let
\vspace{-0.1cm}
\[
P=
\begin{bmatrix}
|&|&&|\\
\vect{v}_1&\vect{v}_2&\cdots&\vect{v}_n\\
|&|&&|
\end{bmatrix}
\]

contain linearly independent eigenvectors.

Let
\vspace{-0.1cm}
\[
D=
\begin{bmatrix}
\lambda_1&0&\cdots&0\\
0&\lambda_2&\cdots&0\\
\vdots&\vdots&\ddots&\vdots\\
0&0&\cdots&\lambda_n
\end{bmatrix}.
\]

Then
\vspace{-0.1cm}
\[
\boxed{
A=PDP^{-1}.
}
\]

This is called \textbf{diagonalization}.

\end{frame}


% =========================================================
% WHY DIAGONALIZATION
% =========================================================

\begin{frame}{Why Diagonalization is Useful}

For a diagonalizable matrix,

\[
A=PDP^{-1}.
\]

Therefore,

\[
A^k
=
PD^kP^{-1}.
\]

Since $D$ is diagonal,

\[
D^k
=
\begin{bmatrix}
\lambda_1^k&0&\cdots&0\\
0&\lambda_2^k&\cdots&0\\
\vdots&\vdots&\ddots&\vdots\\
0&0&\cdots&\lambda_n^k
\end{bmatrix}.
\]

Thus, powers of $A$ can be computed through simple powers
of its eigenvalues.

\end{frame}


% =========================================================
% GEOMETRIC INTERPRETATION
% =========================================================

\begin{frame}{Geometric Interpretation}

A general linear transformation may change both the
direction and magnitude of a vector.

Eigenvectors identify special directions for which

\[
\boxed{
A\vect{v}=\lambda\vect{v}.
}
\]

The direction is preserved while the magnitude is scaled
by $|\lambda|$.

\vspace{0.4cm}

Therefore, eigenvectors reveal the transformation's
\textbf{intrinsic directions}.

\end{frame}


% =========================================================
% EIGENVALUE ZERO
% =========================================================

\begin{frame}{Eigenvalue $\lambda=0$}

If
\vspace{-0.1cm}
\[
A\vect{v}=0\vect{v}=\mathbf{0},
\]

for a non-zero $\vect{v}$, then
\vspace{-0.1cm}
\[
\boxed{
\lambda=0
}
\]

is an eigenvalue.

This means that

\[
\vect{v}\in\operatorname{Null}(A).
\]

Therefore,
\vspace{-0.1cm}
\[
\boxed{
0\text{ is an eigenvalue}
\iff
\operatorname{Null}(A)\neq\{\mathbf{0}\}.
}
\]

For a square matrix, this is also equivalent to

\[
\boxed{
\det(A)=0.
}
\]

\end{frame}


% =========================================================
% INVERTIBILITY AND EIGENVALUES
% =========================================================

\begin{frame}{Eigenvalues and Invertibility}

For a square matrix $A$,

\[
A\text{ is invertible}
\]

if and only if

\[
\det(A)\neq0.
\]

Since

\[
\det(A)
=
\lambda_1\lambda_2\cdots\lambda_n,
\]

when the eigenvalues are counted with algebraic
multiplicity,

\[
\boxed{
A\text{ is invertible}
\iff
0\text{ is not an eigenvalue}.
}
\]

Thus, determinants and eigenvalues provide two different
ways to study invertibility.

\end{frame}


% =========================================================
% TRACE
% =========================================================

\begin{frame}{Trace and Eigenvalues}

The \textbf{trace} of a square matrix is the sum of its
diagonal entries:

\[
\boxed{
\operatorname{tr}(A)
=
a_{11}+a_{22}+\cdots+a_{nn}.
}
\]

The trace is equal to the sum of the eigenvalues:

\[
\boxed{
\operatorname{tr}(A)
=
\lambda_1+\lambda_2+\cdots+\lambda_n.
}
\]

Similarly,

\[
\boxed{
\det(A)
=
\lambda_1\lambda_2\cdots\lambda_n.
}
\]

These relationships connect eigenvalues with two fundamental
matrix quantities.

\end{frame}


% =========================================================
% SYMMETRIC MATRICES
% =========================================================

\begin{frame}{Eigenvectors of Symmetric Matrices}

A real matrix is symmetric if

\[
\boxed{
A=A^T.
}
\]

For real symmetric matrices:

\begin{itemize}
    \item all eigenvalues are real,
    \item eigenvectors corresponding to distinct eigenvalues
    are orthogonal,
    \item an orthonormal basis of eigenvectors exists.
\end{itemize}

Therefore, a real symmetric matrix can be written as

\[
\boxed{
A=QDQ^T,
}
\]

where $Q$ is an orthogonal matrix.

\end{frame}


% =========================================================
% EIGENVALUES IN COMPUTING
% =========================================================

\begin{frame}{Applications of Eigenvalues}

Eigenvalues and eigenvectors appear throughout computing.

\begin{itemize}
    \item \textbf{Principal Component Analysis (PCA)}
    \item \textbf{Dimensionality reduction}
    \item \textbf{Image processing}
    \item \textbf{Graph and network analysis}
    \item \textbf{Differential equations}
    \item \textbf{Stability analysis}
\end{itemize}

\vspace{0.4cm}

The central idea is to identify directions in which a
transformation acts by simple scaling.

\end{frame}